%! Operations with Polynomials — Unit 4, Lesson 4.1 — Homework (Answer Key)
\documentclass[10pt]{article}
\usepackage{algebra2-article}
\usepackage{algebra2-key}   % pulls in -boxes; do NOT also load -boxes
\newcommand{\TallMath}[1]{$\displaystyle #1\rule[-1.4em]{0pt}{3.2em}$}

\begin{document}

\pageheader{Unit 4, Lesson 4.1}{Homework --- Answer Key}
\namedateperiod

\vspace{0.10in}

\begin{notesbox}{Practice}
Write every answer in \textbf{standard form}. Show the products or the sign changes you used.

\begin{enumerate}[leftmargin=1.9em, itemsep=11pt]
  \item \textbf{Name it.} Consider $8 - 3x^2 + x^5 - 6x$.
    \begin{enumerate}[label=(\alph*), leftmargin=1.8em, itemsep=4pt]
      \item Standard form: \ans{$x^5 - 3x^2 - 6x + 8$}
      \item Degree: \ans{5} \quad Leading coefficient: \ans{1} \quad Number of terms:
            \ans{4}
      \item Based on the degree and leading coefficient, both arms of its graph: \ans{left down, right up}
    \end{enumerate}

  \item \textbf{Add and subtract.}
    \begin{enumerate}[label=(\alph*), leftmargin=1.8em, itemsep=6pt]
      \item $(6x^3 - x + 4) + (2x^2 + 5x - 9) = $ \ans{$6x^3 + 2x^2 + 4x - 5$}
      \item $(9x^2 - 3x + 2) - (4x^2 - 3x - 6) = $ \ans{$5x^2 + 8$}
    \end{enumerate}

  \item \textbf{Multiply.}
    \begin{enumerate}[label=(\alph*), leftmargin=1.8em, itemsep=6pt]
      \item $-3x^2(2x^3 - x + 5) = $ \ans{$-6x^5 + 3x^3 - 15x^2$}
      \item $(4x - 1)(x + 7) = $ \ans{$4x^2 + 27x - 7$}
      \item $(x - 2)(x^2 + 3x - 5) = $ \ans{$x^3 + x^2 - 11x + 10$}
    \end{enumerate}

  \item \textbf{Special products.} Use the identities --- no box needed.

    \vspace{3pt}
    (a) $(x + 9)^2 = $ \ans{$x^2 + 18x + 81$} \qquad (b) $(2x - 5)^2 = $ \ans{$4x^2 - 20x + 25$}

    \vspace{4pt}
    (c) $(3x + 8)(3x - 8) = $ \ans{$9x^2 - 64$} \qquad (d) $(x + 6y)(x - 6y) = $ \ans{$x^2 - 36y^2$}

  \item \textbf{Back to a graph you already read.} On yesterday's Exit Ticket you read the graph of
        $k(x) = (x+1)^2(x-2)$.
    \begin{enumerate}[label=(\alph*), leftmargin=1.8em, itemsep=5pt]
      \item Expand $(x+1)^2$: \ans{$x^2 + 2x + 1$}
      \item Now multiply by $(x-2)$ and write $k$ in standard form: \ans{$x^3 - 3x - 2$}
      \item Degree: \ans{3} \quad Leading coefficient: \ans{1} \quad End behavior:
            \ans{left down, right up}
      \item Substituting $x=0$ into your standard form gives the $y$-intercept \ans{$(0,-2)$}. Does it
            agree with the graph you read yesterday? \ans{Yes}
    \end{enumerate}

  \item \textbf{Explain.} Why does subtracting a polynomial require changing the sign of \emph{every}
        term inside the parentheses, and not just the first one? Support your answer with a short
        example that goes wrong if you change only the first sign.
        \ansline{The minus sign in front of the parentheses is a factor of $-1$, and the distributive property multiplies that $-1$ by \emph{each} term inside --- not only the first. Example: $(x^2)-(x^2-5)$ should be $x^2-x^2+5=5$. Changing only the first sign gives $x^2-x^2-5=-5$, the wrong answer by $10$.}
\end{enumerate}
\end{notesbox}

\begin{extensionbox}
\textbf{An area model that needs both operations.} A rectangular garden measures $(2x+3)$ feet by
$(x+5)$ feet. A square pond with side $x$ feet is dug inside it, and the rest is planted.
\begin{enumerate}[label=(\alph*), leftmargin=1.8em, itemsep=5pt]
  \item Write and expand a polynomial for the \emph{total} garden area.
        \ans{$(2x+3)(x+5) = 2x^2 + 13x + 15$ square feet.}
  \item Write a polynomial for the \emph{planted} area (total minus pond) in standard form.
        \ans{$(2x^2+13x+15) - x^2 = x^2 + 13x + 15$ square feet.}
  \item Use your answer to (b) to find the planted area when $x = 4$ feet.
        \ans{$16 + 52 + 15 = 83$ square feet.}
  \item Check part (c) a second way: find the garden's actual dimensions and the pond's actual area at
        $x=4$, and subtract. Do the two methods agree?
        \ansline{At $x=4$ the garden is $11$ ft by $9$ ft, so its area is $99$ ft\textsuperscript{2}; the pond is $4\times4=16$ ft\textsuperscript{2}. Then $99-16=83$ ft\textsuperscript{2} --- the same as (c), so the polynomial is correct.}
\end{enumerate}
\end{extensionbox}

\begin{spiralbox}
\textbf{Coming next --- Lesson 4.2: Advanced Factoring.} Today you ran products \emph{forward}:
$(x-1)^2(x+2)$ became $x^3-3x+2$. Next you will run them \emph{backward} --- pulling out a GCF,
factoring by grouping, and recognizing the difference of squares and the new sum and difference of
\emph{cubes} --- turning a standard form back into the factored form that shows you the zeros.
\end{spiralbox}

\begin{teachernote}
Item~2(b) is the diagnostic: the $-3x$ terms cancel and the constants become $+8$, so an answer of
$5x^2-6x-4$ means the subtraction was distributed to no term but the first. Item~5 is the one to
grade for understanding rather than points --- $k(x)=x^3-3x-2$ should match yesterday's graph at
every feature students already recorded (zeros $-1$ and $2$, $y$-intercept $(0,-2)$, left down / right
up), and naming that agreement out loud is what makes ``two forms, one function'' stick. The
extension's part~(d) is the same idea a third time: a polynomial model is only trustworthy if it
agrees with a direct computation.
\end{teachernote}

\end{document}
